We train a conditional GAN to generate MFTR envelope samples conditioned on the MFTR parameters. Each training example is a pair $(x,\mathbf{c})$, where $x\in\mathbb{R}$ is a single amplitude sample and $\mathbf{c}=[m,\mu,K,\Delta,\bar{\gamma}]$ is the conditioning vector.

\subsection{Architecture}
The model has two multilayer perceptrons (MLPs): a generator $G$ and a discriminator $D$.
The generator takes a noise vector $\mathbf{z}\in\mathbb{R}^{1024}$ and the condition $\mathbf{c}$, embeds $\mathbf{c}$ with a small MLP (Linear--ReLU--Linear) to a 32-dimensional vector, concatenates $[\mathbf{z},\mathrm{emb}(\mathbf{c})]$, and passes it through an MLP with hidden sizes 512--256--128 to output one value. All intermediate layers have LeakyReLU with $\beta 0.2$ as activation function, and \texttt{Softplus} activation is used at the output to keep generated amplitudes positive.
The discriminator embeds the condition the same way, concatenates $[x,\mathrm{emb}(\mathbf{c})]$, and uses an MLP with hidden sizes 512--256--128 and the same activation functions as the generator (except the softplus) to output a single logit.

\subsection{Loss function and training tricks}
We use the standard GAN objective implemented with \texttt{BCEWithLogitsLoss} (numerically stable because $D$ outputs logits).
For the discriminator we use
\begin{equation}
\mathcal{L}_D = \mathrm{BCE}(D(x,\mathbf{c}), 0.9) + \mathrm{BCE}(D(G(\mathbf{z},\mathbf{c}),\mathbf{c}), 0),
\end{equation}
where the target for real samples is set to 0.9 (one-sided label smoothing).
For the generator we use the non-saturating loss
\begin{equation}
\mathcal{L}_G = \mathrm{BCE}(D(G(\mathbf{z},\mathbf{c}),\mathbf{c}), 1).
\end{equation}
We update $D$ every step and update $G$ every $k$ steps (here $k=2$). Linear layers are initialized with Xavier uniform initialization.

\subsection{Training setup}
Training is done for 150 epochs using Adam for both networks, with $\beta_1=0.5$ and $\beta_2=0.999$.
We use learning rates $\mathrm{lr}_G=2.8\times 10^{-4}$ and $\mathrm{lr}_D=1.12\times 10^{-3}$ and a batch size of 512.
% During training we monitor a distributional validation metric: the 1D Wasserstein distance between real and generated validation samples (computed on a fixed, seeded subset for comparability, and pooled across devices when using distributed training). This metric is used for checkpointing and early stopping.

\subsection{Condition normalization and embedding}
The five condition parameters are min--max normalized to $[-1,1]$ using the fixed training ranges, independently per dimension. If a parameter is fixed (low$=$high), its normalized value is set to 0.
Both $G$ and $D$ use the same conditioning mechanism: the normalized condition vector is projected to a learned embedding of dimension 32 using the small embedding MLP described above.

\subsection{Dataset generation (uniform combos and varying only $\mu$)}
We construct a conditional dataset by sampling a set of parameter combinations (\emph{combos}) uniformly within the chosen ranges, and then generating many MFTR samples for each combo.
In this report we focus on a simplified experiment where only $\mu$ varies uniformly in $[1,9]$, while the other parameters are fixed ($m=8$, $K=9$, $\Delta=0.9$, $\bar{\gamma}=1$). This setting is useful because changing $\mu$ can move the MFTR envelope PDF from unimodal to bimodal, so it tests whether the model can smoothly adapt its output shape.
We sample 128 training combos and 32 validation combos uniformly (disjoint split). For each combo we draw 10{,}000 training samples (and 1{,}000 validation samples) from the analytic MFTR sampler (amplitude mode). The dataset stores the combos in a tensor of shape (combos, 5) and returns samples with their corresponding condition by repeating each combo for its associated samples.
